\scanpage{5}
que l'effet d'une permutation circulaire des vecteurs de base sur la matrice $\varphi$ s'obtient par une permutation circulaire sur les \inserted{trois groupes de} lettres $(a,b,c)$, $(p,q,r)$, $(p',q',r')$. Ceci posé, les \deleted{\unsure{la}} formules (13) (13 bis) se récrivent
\[
\left\{\begin{aligned}
 [[e'_1\wedge e'_2],e'_3]&=(ap'-qr')e'_1+(p'r-bq)e'_2+(p'q'-pq)e'_3\\
 [[e'_2\wedge e'_3],e'_1]&=(q'r'-qr)e'_1+(bq'-rp')e'_2+(q'p-cr)e'_3\\
 [[e'_3,e'_1],e'_2]&=(r'q-ap)e'_1+(r'p'-rp)e'_2+(cr'-pq')e'_3
\end{aligned}\right. \tag*{(16)}
\]
d'où
\[
\begin{aligned}
 \text{Jác}(e_1,e_2,e_{\unsure{1|3}})&=[[e_1,e_2],e_3]+[[e_2,e_3],e_1]+[[e_3,e_1],e_2]\\
 &=[a(p'-p)+(q'r'-qr)]e'_1+[b(q'-q)+(r'p'-rp)]e'_2+[c(r'-r)+(p'q'-pq)]e'_3
\end{aligned}\tag*{[17]}
\]
Donc, pour que l'Algèbre $E'_\varphi$ ($E$ muni de $[\, ,\,]_\varphi$) soit une Algèbre de Lie, il faut et il suffit qu'elle satisfasse les conditions
\[
\left\{\begin{aligned}
 a(p'-p)+(q'r'-qr)&=0 &\text{ou}\quad (ap'-qr)-(ap-q'r')&=0\\
 b(q'-q)+(r'p'-rp)&=0 &\text{ou}\quad (bq'-rp)-(bq-r'p')&=0\\
 c(r'-r)+(p'q'-pq)&=0 &\text{---}\quad(cr'-pq)-(cr-p'q')&=0
\end{aligned}\right. \tag*{(18)}
\]

\uline{NB} Avec la notation évidente pour les mineurs d'ordre 2 $M_{\substack{\unsure{i|1|k}\\\unsure{j|3|\ell}}}^{\unsure{k\ell|12|ij}}(\varphi)$ d'une matrice $\varphi$, les premiers membres des équations (18), i.e. les coefficients de l'alternateur jacobien (18), sont
\[
\left\{\begin{aligned}
 &\operatorname{Min}_{\substack{1\\3}}^{12}(\varphi)-\operatorname{Min}_{\substack{1\\3}}^{12}({}^{t}\varphi)\\
 &\operatorname{Min}_{\substack{2\\1}}^{23}(\varphi)-\operatorname{Min}_{\substack{2\\1}}^{23}({}^{t}\varphi)\\
 &\operatorname{Min}_{\substack{3\\2}}^{31}(\varphi)-\operatorname{Min}_{\substack{3\\2}}^{\unsure{31|21}}({}^{t}\varphi)
\end{aligned}\right. \tag*{(19)}
\]

% archive sheet number: 4
